% !TeX root = ../master.tex

\chapter{Introduction}
\label{ch:introduction}

% Bullet points / ideas:
% - Feger et al. (ACL 2025): BERT/RoBERTa/DistilBERT learn datasets, not arguments
%   - Strong in-distribution (F1=0.79), terrible cross-dataset transfer (F1=0.56--0.61)
%   - Manipulation experiments: removing stop words/function words/discourse markers -> delta <= 0.02
%   - Models rely on content-word shortcuts, not argument structure
% - Their conclusion: "decoder-based argument mining may be of interest" -> this thesis
% - GAIC shared task (Touche @ CLEF 2026) provides infrastructure to test this
% - Unique opportunity: do LLMs overcome shortcut learning? do they use context?

\section{Problem Statement}
\label{sec:problem_statement}

% - Current AM systems generalize poorly across datasets
% - Shortcut learning as root cause (Geirhos et al.)
% - Open question: is this architecture-dependent (encoders) or training-dependent?
% - Zero-shot LLMs as natural test case: no dataset-specific training -> no learned shortcuts?

\section{Research Questions}
\label{sec:research_questions}

% RQ1: Do zero-shot LLMs rely on argument structure or superficial cues?
%   - Hypothesis: larger delta than encoders under Feger manipulation + shuffle
% RQ2: Does rich context (definitions, guidelines, document context) improve identification?
%   - Hypothesis: annotation guidelines improve F1, especially dataset-specific ones
% RQ3: Does fine-tuning reintroduce shortcut learning patterns seen in encoders?
%   - Hypothesis: fine-tuning reduces delta, indicating shortcut emergence

\section{Contributions}
\label{sec:contributions}

% - First systematic encoder vs. LLM shortcut learning comparison in AM
% - Word-shuffle manipulation as complementary diagnostic to Feger et al.
% - Quantification of context utilization across 10 datasets and 5 models
% - Evidence on architecture-dependent vs. training-dependent shortcut learning
% - Competitive GAIC 2026 submission + reproducible codebase

\section{Thesis Structure}
\label{sec:thesis_structure}

% - Ch 2: Background (AM, shortcut learning, LLMs, related work)
% - Ch 3: GAIC task description
% - Ch 4: Methodology (shared across all parts)
% - Ch 5: Part 1 -- Zero-shot robustness experiments
% - Ch 6: Part 2 -- Context utilization experiments
% - Ch 7: Part 3 -- Fine-tuning effects
% - Ch 8: Discussion (integration, comparison to encoders, limitations)
% - Ch 9: Conclusion
